% =====================================================================
\chapter{Autoevaluación}
\label{ap:autoevaluacion}
% =====================================================================

\section{Cálculo de la nota}

La nota del examen es la suma de los puntos de la parte~A (máximo 30) y de la parte~B (máximo 70). El examen se aprueba con \textbf{70 puntos o más}. Si se obtienen menos de 18 puntos en la parte~A o menos de 45 en la parte~B, conviene repasar esa parte aunque el examen esté aprobado: el proyecto de fase usa todo este material.

\begin{center}
\begin{tabular}{@{}llcc@{}}
\toprule
Parte & Elemento & Puntos & Obtenidos \\
\midrule
A & Cuestiones 1--20 (1,5 cada una) & 30 & \\
\midrule
B & B.1 SFT: máscara, normalización y empaquetado & 10 & \\
B & B.2 LoRA, QLoRA y olvido & 10 & \\
B & B.3 Valor, gradiente de política y línea base & 10 & \\
B & B.4 Preferencias, política óptima y DPO & 10 & \\
B & B.5 Recuperación léxica y evaluación de rankings & 10 & \\
B & B.6 RAG, agentes y evaluación con incertidumbre & 10 & \\
B & B.7 Superposición, SAE y seguridad & 10 & \\
\midrule
 & \textbf{Total} & \textbf{100} & \\
\bottomrule
\end{tabular}
\end{center}

\section{Puntuación por módulo}

La tabla siguiente agrupa los puntos por módulo. Un módulo en el que se obtenga menos del $70\,\%$ de sus puntos debe repasarse aunque el examen esté aprobado.

\begin{center}
\begin{tabular}{@{}lccccc@{}}
\toprule
Módulo & Cuestiones & Puntos A & Ejercicios & Puntos B & Total \\
\midrule
M20 & 1--5 & 7,5 & B.1, B.2 & 20 & 27,5 \\
M21 & 6--11 & 9 & B.3, B.4 & 20 & 29 \\
M22 & 12--16 & 7,5 & B.5, B.6 & 20 & 27,5 \\
M23 & 17--20 & 6 & B.7 & 10 & 16 \\
\midrule
\textbf{Total} & 20 & \textbf{30} & 7 & \textbf{70} & \textbf{100} \\
\bottomrule
\end{tabular}
\end{center}

El apartado B.6(d) usa también el intervalo conservador del M23, §6.1, y se computa en el M22.

\section{Puerta de salida de la Fase V}

Para dar por superada la Fase~V, y con ella el curso, deben cumplirse las tres condiciones siguientes:
\begin{enumerate}[label=\arabic*.]
  \item un resultado de al menos 9 sobre 12 en el test de cada módulo, M20 a M23;
  \item una autoevaluación del laboratorio interactivo de la fase de al menos el $70\,\%$;
  \item una nota de este examen de hito de al menos el $70\,\%$ (70 puntos sobre 100).
\end{enumerate}
El proyecto de fase se califica con la rúbrica de la \cref{cap:proyecto} y cuenta en la nota de la fase.

\section{Qué repasar según los fallos}

\begin{center}
\small
\renewcommand{\arraystretch}{1.1}
\begin{tabularx}{\textwidth}{@{}>{\raggedright\arraybackslash}p{3.6cm}X@{}}
\toprule
Fallo & Secciones de los apuntes \\
\midrule
Cuestión 1 y B.1 & M20, §2.2 (pérdida enmascarada), §2.3 (lotes, empaquetado y acumulación) y §2.4 (plantilla de chat) \\
Cuestión 4 y B.2 (d) & M20, §3.2 (modelo cuadrático del olvido) \\
Cuestiones 2 y 3 y B.2 (a) & M20, §5.1 (parámetros y fusión) y §5.3 (gradientes e inicialización) \\
B.2 (b) & M20, §5.5 (relación con la SVD); M03 (Eckart--Young) \\
Cuestión 5 y B.2 (c) & M20, §6.1, §6.3 y §6.4 (cuantización por bloques, doble cuantización y QLoRA) \\
Cuestión 6 y B.3 (a) & M21, §1.2 (retorno) y §1.3 (funciones de valor y ecuación de Bellman) \\
Cuestión 7 y B.3 (b) y (c) & M21, §2.3 (REINFORCE) y §2.5 (línea base y varianza) \\
Cuestión 8 & M21, §3.2 (objetivo recortado de PPO) \\
B.4 (b) & M21, §3.3 y §3.4 (penalización KL y política óptima) \\
Cuestión 9 y B.4 (a) & M21, §4.1 y §4.2 (Bradley--Terry y máxima verosimilitud) \\
Cuestión 10 y B.4 (c) y (d) & M21, §5.1 y §5.3 (reparametrización y pérdida DPO) \\
Cuestión 11 y B.4 (e) & M21, §6.3 (GRPO) \\
Cuestión 12 y B.5 (a) & M22, §2.2 y §2.3 (tf-idf y espacio vectorial) \\
Cuestión 13 y B.5 (b) & M22, §2.6 (BM25) \\
Cuestión 14 y B.5 (c) y (d) & M22, §3.2 (métricas de un ranking) \\
Cuestión 15 y B.6 (d) & M22, §3.3 y §3.4 (error estándar y contaminación); M23, §6.1 \\
B.6 (a) y (b) & M22, §4.2 (coseno y producto escalar) y §5.3 (troceo) \\
Cuestión 16 y B.6 (c) & M22, §6.4 (fiabilidad de una cadena de pasos) \\
Cuestión 17 y B.7 (a) & M23, §3.1, §3.3 y §3.4 (superposición, modelo de juguete y dimensión por característica) \\
Cuestión 18 y B.7 (b) & M23, §4.2 y §4.3 (autoencoder disperso y umbral suave) \\
Cuestión 19 y B.7 (c) & M23, §6.3 y §6.4 (calibración y reglas de puntuación) \\
Cuestión 20 y B.7 (d) & M23, §6.6 (estadística de los sucesos raros) \\
\bottomrule
\end{tabularx}
\end{center}

\section{Qué hacer si no se alcanza el umbral}

\begin{enumerate}[label=\arabic*.]
  \item Anotar cada fallo en el registro de errores, clasificado como error de \emph{concepto}, de \emph{cálculo} o \emph{despiste}.
  \item Para cada módulo por debajo del $70\,\%$, repetir el ciclo de estudio del método del curso con las secciones indicadas en la tabla anterior: teoría, laboratorio, práctica, ejercicios $\star\star\star$ del tema, test y tarjetas de repaso.
  \item Rehacer los ejercicios fallados sin mirar la solución y, a continuación, los ejercicios del mismo tema en los apuntes del módulo y las tareas universitarias citadas en la bibliografía.
  \item Repetir el examen pasada al menos una semana, preferiblemente con datos cambiados, y no dar por cerrada la fase hasta superar las tres condiciones de la puerta de salida.
\end{enumerate}
